\documentclass[11pt]{amsart}

\usepackage[T1]{fontenc}
\usepackage[utf8]{inputenc}
\usepackage{amsmath,amssymb,amsthm}
\usepackage[margin=1in]{geometry}
\usepackage{hyperref}

\newcommand{\OO}{\mathbb{O}}
\newcommand{\RR}{\mathbb{R}}

\title[Update brief II: the Kirmse twist]{Update brief~II for\\
\emph{An order on the Leech lattice from a $\mathbb{Z}_3$-symmetric
triple-octonion product}\\
\large Historical accuracy of the Kirmse--Coxeter correction}
\author{Jens K\"oplinger}
\date{17 May 2026}

\begin{document}
\maketitle

\begin{abstract}
The manuscript version of 29 April 2026 (v3) is frozen while under review.
This is the \emph{second}, separate update brief for that frozen version;
it is not to be merged with the first
(\texttt{update\_brief\_2026-04-29\_v3.tex}, which covers a footnote
correction and an appendix simplification).  The present brief concerns one
topic only: how the manuscript describes the Kirmse--Coxeter correction of
the integral octonions (the ``Kirmse twist'') --- in the Introduction, in the
two relevant Remarks of Sections~2 and~3, and in the Section~6 historical
paragraph.  \textbf{Its contents are provisional.}  They are to be finalised
only after the primary sources are obtained and the history is re-derived
from scratch.  Detailed reasoning and verbatim source quotations are in
\texttt{paper/reviews/2026-05-17\_claude\_opus\_4\_7\_kirmse\_history\_note.md}.
\end{abstract}

% ===============================================================
\section{Plan}
% ===============================================================

The work proceeds in two steps, in order:

\begin{enumerate}
\item[(1)] \textbf{Verify the history from primary sources.}  Obtain the
  original papers (Section~\ref{sec:sources}), read them, and derive the
  historical account from scratch --- confirming what is correct and
  identifying any errors, in the manuscript and in the project's own
  reference files.  No secondary source (Petersson, Conway--Smith,
  Wikipedia, Baez) is to be taken on trust; they serve only as
  cross-checks.
\item[(2)] \textbf{Elevate the wording.}  Once the history is settled,
  rewrite the description of the Kirmse twist.  The current wording ---
  ``swapping two basis indices in the Fano-plane labelling'', ``there are
  seven of them'' --- is imprecise and is to be replaced by an accurate,
  structurally explicit account.
\end{enumerate}

\noindent
Everything in Sections~\ref{sec:sources}--\ref{sec:items} below is
\emph{pending} step~(1).  Nothing here is final.

% ===============================================================
\section{Primary sources to obtain}\label{sec:sources}
% ===============================================================

The mathematical facts in this area ($E_8$ is a maximal order; there are
seven; the twist yields a closed order) are to be verified
\emph{computationally} in-project, not cited.  Only the \emph{historical}
record requires the original documents:

\begin{itemize}
\item \textbf{Kirmse (1924)} --- \emph{required.}  J.~Kirmse, ``\"Uber die
  Darstellbarkeit nat\"urlicher ganzer Zahlen als Summen von acht Quadraten
  \ldots,'' Ber.\ Verh.\ S\"achs.\ Akad.\ Wiss.\ Leipzig, Math.-Phys.\ Kl.\
  \textbf{76} (1924), 63--82.  German.  Settles whether Kirmse proposed one
  lattice or seven.
\item \textbf{Coxeter (1946)} --- \emph{required.}  H.\,S.\,M.~Coxeter,
  ``Integral Cayley numbers,'' Duke Math.\ J.\ \textbf{13} (1946), 561--578.
  Settles the precise form of the correction and how Bruck is credited.
  Paywalled online; needs a library.
\item \textbf{Dickson (1923)} --- \emph{recommended; already free online}
  (Numdam).  L.\,E.~Dickson, ``A new simple theory of hypercomplex
  integers,'' J.~Math.\ Pures Appl.\ (9) \textbf{2} (1923), 281--326.
\item \textbf{Mahler (1942)} --- \emph{optional, not load-bearing.}
  K.~Mahler, ``On ideals in the Cayley--Dickson algebra,'' Proc.\ Roy.\
  Irish Acad.\ Sect.\ A \textbf{48} (1942), 123--133.
\end{itemize}

\noindent
There is no separate Bruck publication; his correction is recorded within
Coxeter (1946).

% ===============================================================
\section{Items provisionally identified}\label{sec:items}
% ===============================================================

The following are the candidate changes, \textbf{each pending verification}
in step~(1).  They are recorded here so the digest has a starting point.

\begin{itemize}
\item \textbf{Credit Bruck.}  The secondary record (Petersson 2018)
  attributes the actual correction to R.\,H.~Bruck, with Coxeter publishing
  it.  The manuscript credits only Coxeter.  To be confirmed against
  Coxeter (1946) and then corrected in the prose.
\item \textbf{The correcting operation.}  The secondary record describes it
  as interchanging the identity element with one imaginary unit (seven
  choices $\to$ seven maximal orders), \emph{not} as ``swapping two basis
  indices in the Fano-plane labelling.''  The accurate description --- and
  its precise relation to the manuscript's transposition~$\sigma$ of two
  \emph{imaginary} units --- is to be established from Coxeter (1946).
\item \textbf{Re-examine the Remark contrasting $\sigma$ with a Fano-line
  permutation.}  It currently asserts that Coxeter's correction is ``not
  induced by a linear map on~$\RR^8$.''  If the correction is a coordinate
  interchange, that assertion is wrong and the intended contrast must be
  restated.  To be settled from Coxeter (1946).
\item \textbf{Did Kirmse propose seven orders?}  The secondary sources
  disagree.  This is the central open historical question and is to be
  decided from Kirmse (1924) itself, cross-checked with Coxeter (1946).
\item \textbf{Dickson's priority.}  Dickson (1923) reportedly gave a correct
  integral-octonion construction more than twenty years before Coxeter.
  Whether, and how, to record this (and to cite Dickson, and optionally
  Mahler) is to be decided after reading Dickson (1923).
\item \textbf{Project reference files.}  \texttt{background\_octonions.txt}
  and \texttt{terminology.md} are to be checked for the same issues,
  including a discrepancy in the name of Kirmse's journal.
\end{itemize}

% ===============================================================
\section*{Note on the companion document}
% ===============================================================

\texttt{companion.tex} carries the same imprecise wording.  It is a frozen
training document and is not edited here, but the same correction should be
applied if the companion is ever revised.

\end{document}
